% Part: lambda-calculus
% Chapter: lambda-definability
% Section: lists

\documentclass[../../../include/open-logic-section]{subfiles}

\begin{document}

\olfileid{lam}{ldf}{lst}
\olsection{યાદીઓ}

યાદી $[M_1, M_2, \ldots, M_n]$ આ પ્રમાણે
વ્યાખ્યાયિત કરી શકાય છે:
\[
  [M_1, M_2, \ldots, M_n] \ident \lambd[sz][s M_1 (s M_2 (\ldots (s M_n z)))]
\]
અનૌપચારિક રીતે, યાદીનું નિરૂપણ કરતું પદ
તે યાદીનું ``સંચયક'' વિધેય છે. તે બે પદો
લે છે. પહેલું પદ નવો ઘટક અને અત્યાર સુધીનું
સંચિત મૂલ્ય લઈને નવું સંચિત મૂલ્ય આપતું
વિધેય છે; બીજું પદ શરૂઆતનું સંચિત મૂલ્ય છે.
આ રીતે તે એક પછી એક ઘટક સંચિત કરીને અંતે
પરિણામ આપે છે.
\sourcecorrection{OLMD-187}{મૂળ ગદ્યમાં
પહેલા વિધેયને સંચિત મૂલ્ય અને નવો ઘટક
એવા ક્રમમાં આર્ગ્યુમેન્ટ મળતાં કહે છે.
પ્રદર્શિત પદ $s\,M_i\,a$માં નવો ઘટક
પહેલો અને સંચિત મૂલ્ય બીજું છે; અહીં
ગદ્યનો ક્રમ પદ મુજબ રાખ્યો છે.}

\begin{digress}
  વિધેયાત્મક પ્રોગ્રામિંગથી પરિચિત હો
  તો આ વ્યાખ્યા \emph{જમણી ફોલ્ડ} જેવી
  લાગશે: યાદી તેના ઘટકોના જમણી ફોલ્ડ
  વિધેય તરીકે વ્યાખ્યાયિત થાય છે.
\end{digress}

આ સંકેતનના આધારે કેટલાંક ઉપયોગી વિધેયો
સરળતાથી વ્યાખ્યાયિત કરી શકાય છે:
\begin{align*}
  \fn{Sum} & \ident
  \lambd[l][l \, (\lambd[xa][\fn{Add} \, x \, a])\,  \num{0}]\\
  \fn{Len} & \ident
  \lambd[l][l \, (\lambd[xa][\fn{Add} \, a \, \num{1}]) \num{0}]
\end{align*}

$\fn{Sum}$ ચર્ચ અંકોની યાદીનો સરવાળો
કાઢે છે. યાદી પર સંચયન કરતાં તેનું
શરૂઆતનું મૂલ્ય $\num{0}$ છે અને દરેક
ઘટક માટે નવું મૂલ્ય $\fn{Add} \, x\, a$
છે; અહીં $x$ હાલનો ઘટક અને $a$ સંચિત
મૂલ્ય છે. આખરે બધા ઘટકોનો સરવાળો મળે છે.

\begin{prob}
  $\fn{Len}$ શું કરે છે? સમજાવો.
\end{prob}

\begin{prob}
  યાદીનો ક્રમ ઊલટાવતું વિધેય વ્યાખ્યાયિત કરો.
\end{prob}
\end{document}
